% !TEX program = xelatex
% 俄语 · 形容词总表（按课件分类）· 2026-11-13
\documentclass[11pt]{ctexart}
%=====================================================================
% theme-violet.tex —— 讲义精讲 · 紫花主题
% 用法：\documentclass[11pt]{ctexart}
%       \input{assets/theme-violet}
%       \begin{document} ... \end{document}
% 只用 ctexart 自带能力 + xcolor/enumitem/fancyhdr/pifont/amsmath，
% 不依赖 tikz/tcolorbox，保证 MiKTeX 可直接编译。
%=====================================================================

\usepackage[a4paper,left=2.0cm,right=2.0cm,top=2.1cm,bottom=2.3cm,%
            headsep=12pt,footskip=22pt]{geometry}
\usepackage{amsmath,amssymb}
\usepackage{needspace}
\usepackage{xcolor}
\usepackage{array}
\usepackage{enumitem}
\usepackage{pifont}
\usepackage{fancyhdr}
\usepackage[hidelinks]{hyperref}

%---------------------------------------------------------------- 配色
\definecolor{violet}{RGB}{107,78,143}      % 主色：紫
\definecolor{violetdeep}{RGB}{63,44,92}    % 深紫：标题/强调
\definecolor{petal}{RGB}{183,159,208}      % 花瓣：装饰线、花饰
\definecolor{bandviolet}{RGB}{247,243,251} % 浅紫底：档案/提示框
\definecolor{bandcream}{RGB}{252,249,242}  % 米白底：易错/注意框
\definecolor{rulepurple}{RGB}{217,207,230} % 分隔线
\definecolor{inkgray}{RGB}{92,96,110}      % 次要文字

%---------------------------------------------------------------- 字体
% 有 Noto CJK 用 Noto，没有就退回 ctex 默认（Windows 下为宋体/黑体）
\IfFontExistsTF{Noto Serif CJK SC}{\setCJKmainfont{Noto Serif CJK SC}}{}
\IfFontExistsTF{Noto Sans CJK SC}{\setCJKsansfont{Noto Sans CJK SC}}{}

%---------------------------------------------------------------- 版式
\setlist{nosep,leftmargin=1.9em}
\linespread{1.28}
\setlength{\parskip}{3.5pt}
\setlength{\parindent}{0pt}
\setlength{\tabcolsep}{5pt}
\renewcommand{\arraystretch}{1.28}

%---------------------------------------------------------------- 花饰
% \ding{95} = 五瓣小花，\ding{96} = 星形花；都来自 pifont，必定存在
\newcommand{\orn}{\textcolor{petal}{\ding{95}}}          % 行内小标
\newcommand{\ornstar}{\textcolor{petal}{\ding{96}}}      % 星形花
\newcommand{\flowerrule}{%                                % 居中花饰分隔线
  \par\vspace{7pt}%
  \noindent\makebox[\linewidth]{\color{petal}%
    \rule{0.20\linewidth}{0.6pt}\hspace{8pt}{\small\ding{95}}\hspace{4pt}%
    \raisebox{1.1pt}{\tiny\ding{95}}\hspace{4pt}{\small\ding{95}}\hspace{8pt}%
    \rule{0.20\linewidth}{0.6pt}}%
  \par\vspace{7pt}}

%---------------------------------------------------------------- 标题与节头
\newcommand{\doctitle}[3]{%   #1 主标题  #2 副标题  #3 元信息（可用 \\ 分行）
  \begin{center}
    {\color{petal}\small \orn\enspace #2\enspace \ornstar}\\[7pt]
    {\LARGE\bfseries\color{violetdeep} #1}\\[10pt]
    {\color{petal}\rule{0.55\linewidth}{1pt}}\\[4pt]
    {\small\color{inkgray} #3}
  \end{center}\vspace{2pt}}

\newcommand{\hdr}[2]{%   #1 章节号（如 §0）  #2 标题
  \par\vspace{12pt}%
  \noindent\colorbox{violetdeep}{\color{white}\bfseries\small\;#1\;}\hspace{7pt}%
  {\color{violetdeep}\large\bfseries #2}\par\nobreak
  \vspace{3pt}{\color{petal}\hrule height 1.2pt}\vspace{6pt}}

\newcommand{\subhdr}[1]{%
  \par\vspace{6pt}%
  \noindent{\color{violet}\bfseries #1}\par\vspace{1pt}}

%---------------------------------------------------------------- 提示框
\newcommand{\callout}[1]{% 浅紫底：讲义档案、读法说明
  \par\medskip\noindent
  \fcolorbox{rulepurple}{bandviolet}{\parbox{\dimexpr\linewidth-2\fboxsep-2\fboxrule}{\vspace{2pt}#1\vspace{2pt}}}%
  \par\medskip}

\newcommand{\warnbox}[1]{% 米白底：易错、边界条件
  \par\medskip\noindent
  \fcolorbox{rulepurple}{bandcream}{\parbox{\dimexpr\linewidth-2\fboxsep-2\fboxrule}{\vspace{2pt}#1\vspace{2pt}}}%
  \par\medskip}

\newcommand{\gapitem}[2]{% #1 缺口类型（紫） #2 内容
  \par\vspace{2pt}\noindent
  \hangindent=1.4em\hangafter=1
  {\color{petal}\small\ding{96}}\enspace{\bfseries\color{violetdeep}#1}\enspace #2\par}

% 题干块（\qstem）：每道例题／作业的「详解」开头必须先复述题干与出处。
% 起因：2026-10-25 第五章讲义把七道作业的题面集中放在一页、详解另起一页，
% 读者翻到「第 4 题」时看不到条件，等于一道没有题干的题目。
% 规则：详解块内自带题干，读者不翻页就能读懂题在问什么。
\newcommand{\qstem}[1]{%
  \par\medskip\noindent
  \fcolorbox{violet}{bandviolet}{%
    \parbox{\dimexpr\linewidth-2\fboxsep-2\fboxrule}{%
      \vspace{2pt}%
      {\small\bfseries\color{violet}题干}\hspace{0.8em}#1%
      \vspace{2pt}}}%
  \par\medskip}

% v2: 保护题干与开头，长题可跨页，续页补上下文。
\newcommand{\qbreak}{\needspace{6\baselineskip}}
% 独立练习：ID、来源、题干、考点、答案标签（通常 ans:ID）。
\newcommand{\practice}[5]{%
  \qbreak\subhdr{练习 #1}\qstem{（#2）#3}%
  \textbf{考点}：#4\par
  详解见附录第~\pageref{#5}~页。\par\medskip}
% 附录开头：ID、来源、条件摘要；自动产生 ans:ID 标签。
\newcommand{\answerstart}[3]{%
  \qbreak\subhdr{#1\quad 详解}\label{ans:#1}\qstem{（#2）#3}}
\newcommand{\answercontinue}[2]{%
  \needspace{5\baselineskip}\subhdr{#1\quad 续解}%
  {\small\color{inkgray}本页接续：#2}\par\medskip}
% 延迟提供默认值，让旧讲义可在导言区保留自己的 \newcommand{\step}。
\AtBeginDocument{\providecommand{\step}[1]{\needspace{3\baselineskip}\par\medskip\textbf{#1}\par}}

% 前提旁注：只放在「需要理解的那一步」后面，给陈述 + 本题代入，不重证。
\newcommand{\aside}[1]{%
  \par\vspace{1pt}\noindent
  \hangindent=1.6em\hangafter=1
  {\small\color{inkgray}%
    \hspace*{0.6em}{\color{petal}$\triangleright$}\hspace{2pt}#1}%
  \par\vspace{2pt}}

%---------------------------------------------------------------- 表格
% 本机 MiKTeX 的 colortbl 与 array 版本不匹配（\rowcolor 会报错），
% 故表格统一用「上粗紫线 + 表头 + 中细线」的三线式，不依赖 colortbl。
\newcommand{\tabtop}{\par\vspace{3pt}{\color{violet}\hrule height 1.2pt}\vspace{1pt}}
\newcommand{\tabmid}{\par\vspace{1pt}{\color{petal}\hrule height 0.5pt}\vspace{1pt}}
\newcommand{\tabbot}{\par\vspace{3pt}{\color{rulepurple}\hrule height 0.5pt}\vspace{8pt}}
\newcommand{\thead}[1]{{\bfseries\color{violetdeep}#1}}

%---------------------------------------------------------------- 页眉页脚
\pagestyle{fancy}
\fancyhf{}
\renewcommand{\headrulewidth}{0.4pt}
\renewcommand{\headrule}{\hbox to\headwidth{\color{rulepurple}\leaders\hrule height \headrulewidth\hfill}}
\fancyhead[L]{\small\color{inkgray}\songti \doctag}
\fancyhead[R]{\small\color{inkgray}\songti \coursetag}
\fancyfoot[C]{\small\color{inkgray}\orn\enspace{\color{violet}\thepage}\enspace\orn}

% 每篇讲义在导言区设这两个值
\newcommand{\doctag}{}
\newcommand{\coursetag}{}
\newcommand{\settags}[2]{\renewcommand{\doctag}{#1}\renewcommand{\coursetag}{#2}}

%---------------------------------------------------------------- 数学记号
\newcommand{\vect}[1]{\boldsymbol{#1}}
\newcommand{\dd}{\,\mathrm{d}}
\newcommand{\ddx}{\frac{\mathrm{d}}{\mathrm{d}x}}
\addtolength{\jot}{2pt}



\setmainfont{Times New Roman}
\setsansfont{Times New Roman}
\usepackage{longtable}
\setlength{\headheight}{22pt}
\settags{俄语 · 形容词总表}{俄语（第二外语）}
\newcommand{\blank}{\underline{\hspace{3.2em}}}
\begin{document}

\doctitle{俄语形容词总表}
{按课件分类 · 五族词尾 · 中俄双语}
{УРОК 1--8 里出现的形容词全部收进来，按课件顺序分块；\\
每一条给出原形与阴／中／复三个形式，并标出它属于哪一族 \quad 2026-11-13}

\hdr{§0}{怎么用这份表}

\callout{\textbf{形容词只有一个难点}：词尾要跟着名词的性、数变。课件把它压成\textbf{五族}——硬词干、г／к／х、ж／ш／ч／щ、软词干、词尾重音。所以这份表按「族」组织：每条的\textbf{族}一栏告诉你它走哪一套词尾，右边三格直接给出阴、中、复三个形式。考试里的构形题（Это … (интересный) книга）只要做对两件事：\textbf{认出原形、选对族}。}

\subhdr{0.1　四栏怎么读}

\begin{minipage}{\linewidth}
\tabtop
\begin{tabular}{@{}p{0.14\linewidth}p{0.30\linewidth}p{0.46\linewidth}@{}}
\thead{栏} & \thead{写法} & \thead{意思} \\
\end{tabular}
\tabmid
\begin{tabular}{@{}p{0.14\linewidth}p{0.30\linewidth}p{0.46\linewidth}@{}}
原形 & но́вый & 阳性单数第一格，回答 Како́й?；词典、单词表都只收这一格 \\
中文 & 新的 & 取课件语境里的常用义 \\
族 & 硬／гкх／жшчщ／软／-ой & 词干类型，决定挂哪一套词尾（0.2） \\
三形式 & но́вая · но́вое · но́вые & 阴性（она́）、中性（оно́）、复数（они́） \\
\end{tabular}
\tabbot
\end{minipage}

\subhdr{0.2　五个族：为什么有的写 -ый，有的写 -ий}

\begin{minipage}{\linewidth}\small
\tabtop
\begin{tabular}{@{}p{0.15\linewidth}p{0.20\linewidth}p{0.16\linewidth}p{0.16\linewidth}p{0.16\linewidth}@{}}
\thead{族} & \thead{原形} & \thead{阴性} & \thead{中性} & \thead{复数} \\
\end{tabular}
\tabmid
\begin{tabular}{@{}p{0.15\linewidth}p{0.20\linewidth}p{0.16\linewidth}p{0.16\linewidth}p{0.16\linewidth}@{}}
硬词干 & но́вый & но́вая & но́вое & но́вые \\
г／к／х & высо́кий & высо́кая & высо́кое & высо́кие \\
ж, ш, ч, щ & хоро́ший & хоро́шая & хоро́шее & хоро́шие \\
软词干 & си́ний & си́няя & си́нее & си́ние \\
词尾重音 & голубо́й & голуба́я & голубо́е & голубы́е \\
\end{tabular}
\tabbot
\end{minipage}

\textbf{一句话记法}：\textbf{词干硬就 -ый／-ая／-ое／-ые；词干末是 г、к、х 就写 -ий（阴 -ая、中 -ое、复 -ие）；词干末是 ж、ш、ч、щ 也写 -ий，但中性换 -ее、复数仍是 -ие；词干软就 -ий／-яя／-ее／-ие；重音落在词尾就整族换 -о́й／-а́я／-о́е／-ы́е}。

\warnbox{\textbf{两个反例，考试最爱考}：\textbf{голубы́е}（浅蓝的复数）写 ы，因为词干末尾是 б，不在那七个字母里；\textbf{больши́е}（大的复数）写 и，因为词干末尾是 ш。判断依据永远是\textbf{词干最后一个字母}，不是整个词看起来像谁。}

\subhdr{0.3　建议的三遍读法}

\begin{enumerate}
\item \textbf{第一遍（按课）}：从 §1 读到 §10。每块标题后面的页码就是课件出处，可以回原 PDF 对。
\item \textbf{第二遍（遮右边）}：只看原形，先说中文、再说族、最后念出阴／中／复三个形式。
\item \textbf{第三遍（考前）}：只看 §11 速查卡，加上自己画过道的词。
\end{enumerate}

\begin{minipage}{\linewidth}
\subhdr{0.4　目录}
\tabtop
\begin{tabular}{@{}p{0.09\linewidth}p{0.76\linewidth}p{0.09\linewidth}@{}}
§1 & УРОК 1–6：正式讲形容词之前只出现过两个 & \pageref{sec:1} \\
§2 & УРОК 7（一）：како́й 的四个形式与第一批形容词 & \pageref{sec:2} \\
§3 & УРОК 7（二）：一个形容词 × 四个性数 & \pageref{sec:3} \\
§4 & УРОК 7（三）：13 个颜色词 & \pageref{sec:4} \\
§5 & УРОК 7（四）：季节、关系与顺序 & \pageref{sec:5} \\
§6 & УРОК 7（五）：拼写规则第一组（г／к／х） & \pageref{sec:6} \\
§7 & УРОК 7（六）：拼写规则第二组（ж／ш／ч／щ） & \pageref{sec:7} \\
§8 & УРОК 7（七）：味道与温度 & \pageref{sec:8} \\
§9 & УРОК 7（八）：课文与 нра́виться 练习里的形容词 & \pageref{sec:9} \\
§10 & УРОК 8：са́мый 最高级与指示代词里的形容词 & \pageref{sec:10} \\
§11 & 考前速查：五族、拼写、四个坑 & \pageref{sec:11} \\
§12 & 这份表的来源与边界 & \pageref{sec:12} \\
\end{tabular}
\tabbot
\end{minipage}
\hdr{§1}{УРОК 1–6：正式讲形容词之前只出现过两个}\label{sec:1}
前六课的重心是语音、名词和代词，形容词只在两处露面：У6 p6 的问候语 \textbf{До́брое у́тро! До́брый день! До́брый ве́чер!}，和 У1 p41 的 \textbf{дома́шнее зада́ние}。它们到 У7 才正式讲词尾，这里先认个脸。
\begin{longtable}{@{}p{0.17\linewidth}p{0.19\linewidth}p{0.06\linewidth}p{0.52\linewidth}@{}}
\multicolumn{4}{@{}p{\linewidth}@{}}{\subhdr{1.1　问候与作业里的两个形容词}} \\
\hline
\thead{原形（он）} & \thead{中文} & \thead{族} & \thead{она́／оно́／они́} \\
\hline
\endfirsthead
\hline
\thead{原形（он）} & \thead{中文} & \thead{族} & \thead{она́／оно́／они́} \\
\multicolumn{4}{@{}l}{\scriptsize\color{inkgray}（接上页）1.1　问候与作业里的两个形容词} \\
\hline
\endhead
до́брый & 善良的；问候语里的「好」 & 硬 & {\small до́брая · до́брое · до́брые} \\
дома́шний & 家庭的、家里的 & 软 & {\small дома́шняя · дома́шнее · дома́шние} \\
\hline
\end{longtable}
\hdr{§2}{УРОК 7（一）：како́й 的四个形式与第一批形容词}\label{sec:2}
У7 p26 只讲两件事：形容词永远放在名词前面，词尾跟着名词的性、数变。课件用 \textbf{но́вый телефо́н}、\textbf{краси́вая де́вушка}、\textbf{дли́нная у́лица} 立规矩，p27 再用 \textbf{ста́рое фо́то}、\textbf{у́мные студе́нты} 复一遍。这一批五个词是后面所有练习的种子。
\begin{longtable}{@{}p{0.17\linewidth}p{0.19\linewidth}p{0.06\linewidth}p{0.52\linewidth}@{}}
\multicolumn{4}{@{}p{\linewidth}@{}}{\subhdr{2.1　课件 p26–p30 的五个种子词}\par{\small 课件的四个疑问词直接给出词尾：Како́й? → -ый／-ий／-о́й，Кака́я? → -ая／-яя／-а́я，Како́е? → -ое／-ее／-о́е，Каки́е? → -ые／-ие／-ы́е。答题时照抄问句的词尾即可。}} \\
\hline
\thead{原形（он）} & \thead{中文} & \thead{族} & \thead{она́／оно́／они́} \\
\hline
\endfirsthead
\hline
\thead{原形（он）} & \thead{中文} & \thead{族} & \thead{она́／оно́／они́} \\
\multicolumn{4}{@{}l}{\scriptsize\color{inkgray}（接上页）2.1　课件 p26–p30 的五个种子词} \\
\hline
\endhead
но́вый & 新的 & 硬 & {\small но́вая · но́вое · но́вые} \\
краси́вый & 漂亮的 & 硬 & {\small краси́вая · краси́вое · краси́вые} \\
дли́нный & 长的 & 硬 & {\small дли́нная · дли́нное · дли́нные} \\
ста́рый & 旧的、老的 & 硬 & {\small ста́рая · ста́рое · ста́рые} \\
у́мный & 聪明的 & 硬 & {\small у́мная · у́мное · у́мные} \\
\hline
\end{longtable}
\hdr{§3}{УРОК 7（二）：一个形容词 × 四个性数}\label{sec:3}
p34 的练习把同一个形容词配四个名词：\textbf{краси́вый парк／краси́вая страна́／краси́вое мо́ре／краси́вые де́вушки}，要求当场换词尾。краси́вый、ста́рый 已在 §2 出现过，这一组的新词是下面四个。
\begin{longtable}{@{}p{0.17\linewidth}p{0.19\linewidth}p{0.06\linewidth}p{0.52\linewidth}@{}}
\multicolumn{4}{@{}p{\linewidth}@{}}{\subhdr{3.1　课件 p34 的四组搭配}\par{\small 注意 молодо́й 和 родно́й 是词尾重音族：阳性写 -о́й，阴性 -а́я，复数 -ы́е。весёлый 与 интере́сный 是硬词干，词尾原样贴。}} \\
\hline
\thead{原形（он）} & \thead{中文} & \thead{族} & \thead{она́／оно́／они́} \\
\hline
\endfirsthead
\hline
\thead{原形（он）} & \thead{中文} & \thead{族} & \thead{она́／оно́／они́} \\
\multicolumn{4}{@{}l}{\scriptsize\color{inkgray}（接上页）3.1　课件 p34 的四组搭配} \\
\hline
\endhead
молодо́й & 年轻的 & -ой & {\small молода́я · молодо́е · молоды́е} \\
весёлый & 快乐的 & 硬 & {\small весёлая · весёлое · весёлые} \\
интере́сный & 有趣的 & 硬 & {\small интере́сная · интере́сное · интере́сные} \\
родно́й & 亲的、故乡的 & -ой & {\small родна́я · родно́е · родны́е} \\
\hline
\end{longtable}
\hdr{§4}{УРОК 7（三）：13 个颜色词}\label{sec:4}
颜色是形容词里最整齐的一组，但课件特意把它们分成三种词尾型：11 个走 \textbf{-ый}（含 разноцве́тный），1 个走软词干（\textbf{си́ний}），1 个走词尾重音（\textbf{голубо́й}）。背颜色时把「族」一起背，才不会写出 *си́ная、*голубая 这类错形。
\begin{longtable}{@{}p{0.17\linewidth}p{0.19\linewidth}p{0.06\linewidth}p{0.52\linewidth}@{}}
\multicolumn{4}{@{}p{\linewidth}@{}}{\subhdr{4.1　课件 p36–p37 的颜色表}\par{\small 课件 p37 的顺序是 ора́нжевый、жёлтый、зелёный、голубо́й、си́ний、фиоле́товый、разноцве́тный；这里按「族」排，方便对照词尾。}} \\
\hline
\thead{原形（он）} & \thead{中文} & \thead{族} & \thead{она́／оно́／они́} \\
\hline
\endfirsthead
\hline
\thead{原形（он）} & \thead{中文} & \thead{族} & \thead{она́／оно́／они́} \\
\multicolumn{4}{@{}l}{\scriptsize\color{inkgray}（接上页）4.1　课件 p36–p37 的颜色表} \\
\hline
\endhead
бе́лый & 白色的 & 硬 & {\small бе́лая · бе́лое · бе́лые} \\
чёрный & 黑色的 & 硬 & {\small чёрная · чёрное · чёрные} \\
кра́сный & 红色的 & 硬 & {\small кра́сная · кра́сное · кра́сные} \\
зелёный & 绿色的 & 硬 & {\small зелёная · зелёное · зелёные} \\
жёлтый & 黄色的 & 硬 & {\small жёлтая · жёлтое · жёлтые} \\
се́рый & 灰色的 & 硬 & {\small се́рая · се́рое · се́рые} \\
кори́чневый & 棕色的 & 硬 & {\small кори́чневая · кори́чневое · кори́чневые} \\
ора́нжевый & 橙色的 & 硬 & {\small ора́нжевая · ора́нжевое · ора́нжевые} \\
ро́зовый & 粉色的 & 硬 & {\small ро́зовая · ро́зовое · ро́зовые} \\
фиоле́товый & 紫色的 & 硬 & {\small фиоле́товая · фиоле́товое · фиоле́товые} \\
разноцве́тный & 彩色的 & 硬 & {\footnotesize разноцве́тная · разноцве́тное · разноцве́тные} \\
си́ний & 深蓝色的 & 软 & {\small си́няя · си́нее · си́ние} \\
голубо́й & 浅蓝色的 & -ой & {\small голуба́я · голубо́е · голубы́е} \\
\hline
\end{longtable}
\hdr{§5}{УРОК 7（四）：季节、关系与顺序}\label{sec:5}
p40 的四组练习是 \textbf{весе́нний}、\textbf{осе́нний}、\textbf{дома́шний}、\textbf{после́дний}。дома́шний 在 §1 已出现（дома́шнее зада́ние），这里只补三个。注意这三个都是软词干，阴性和中性写 \textbf{-яя／-ее}。
\begin{longtable}{@{}p{0.17\linewidth}p{0.19\linewidth}p{0.06\linewidth}p{0.52\linewidth}@{}}
\multicolumn{4}{@{}p{\linewidth}@{}}{\subhdr{5.1　课件 p40 的四季与顺序}} \\
\hline
\thead{原形（он）} & \thead{中文} & \thead{族} & \thead{она́／оно́／они́} \\
\hline
\endfirsthead
\hline
\thead{原形（он）} & \thead{中文} & \thead{族} & \thead{она́／оно́／они́} \\
\multicolumn{4}{@{}l}{\scriptsize\color{inkgray}（接上页）5.1　课件 p40 的四季与顺序} \\
\hline
\endhead
весе́нний & 春天的 & 软 & {\small весе́нняя · весе́ннее · весе́нние} \\
осе́нний & 秋天的 & 软 & {\small осе́нняя · осе́ннее · осе́нние} \\
после́дний & 最后的 & 软 & {\small после́дняя · после́днее · после́дние} \\
\hline
\end{longtable}
\hdr{§6}{УРОК 7（五）：拼写规则第一组（г／к／х）}\label{sec:6}
这一组的词干末尾是 \textbf{г、к、х}，阳性只能写 \textbf{-ий}——*ру́скый、*ма́ленькый 都是错写法。дорого́й 虽然也在 p41，但它属于词尾重音那一族，一并放在这里对照。
\begin{longtable}{@{}p{0.17\linewidth}p{0.19\linewidth}p{0.06\linewidth}p{0.52\linewidth}@{}}
\multicolumn{4}{@{}p{\linewidth}@{}}{\subhdr{6.1　课件 p41 的六组练习}\par{\small 课件的模型是 \textbf{плохо́й друг／плоха́я пого́да／плохо́е настрое́ние／плохи́е но́вости}：同一个词拉到四个性数上。注意 плохо́й 是 -о́й 族，复数 плохи́е 写 и（词干末 х）。}} \\
\hline
\thead{原形（он）} & \thead{中文} & \thead{族} & \thead{она́／оно́／они́} \\
\hline
\endfirsthead
\hline
\thead{原形（он）} & \thead{中文} & \thead{族} & \thead{она́／оно́／они́} \\
\multicolumn{4}{@{}l}{\scriptsize\color{inkgray}（接上页）6.1　课件 p41 的六组练习} \\
\hline
\endhead
плохо́й & 坏的 & -ой & {\small плоха́я · плохо́е · плохи́е} \\
ти́хий & 安静的 & гкх & {\small ти́хая · ти́хое · ти́хие} \\
белору́сский & 白俄罗斯的 & гкх & {\small белору́сская · белору́сское · белору́сские} \\
дорого́й & 贵的、亲爱的 & -ой & {\small дорога́я · дорого́е · дороги́е} \\
ру́сский & 俄罗斯的 & гкх & {\small ру́сская · ру́сское · ру́сские} \\
ма́ленький & 小的 & гкх & {\small ма́ленькая · ма́ленькое · ма́ленькие} \\
высо́кий & 高的 & гкх & {\small высо́кая · высо́кое · высо́кие} \\
\hline
\end{longtable}
\hdr{§7}{УРОК 7（六）：拼写规则第二组（ж／ш／ч／щ）}\label{sec:7}
词干末尾是 \textbf{ж、ш、ч、щ}，阳性同样写 \textbf{-ий}，中性写 \textbf{-ее}（хоро́шее，不是 *хоро́шое），复数写 \textbf{-ие}。большо́й、чужо́й 是词尾重音族，但复数仍受这条规则约束：больши́е、чужи́е。
\begin{longtable}{@{}p{0.17\linewidth}p{0.19\linewidth}p{0.06\linewidth}p{0.52\linewidth}@{}}
\multicolumn{4}{@{}p{\linewidth}@{}}{\subhdr{7.1　课件 p42 的六组练习}\par{\small 课件的模型是 \textbf{большо́й го́род／больша́я страна́／большо́е письмо́／больши́е расска́зы}。большо́й 与 чужо́й 的词干以 ш、ж 收尾，所以复数写 и 而不是 ы。}} \\
\hline
\thead{原形（он）} & \thead{中文} & \thead{族} & \thead{она́／оно́／они́} \\
\hline
\endfirsthead
\hline
\thead{原形（он）} & \thead{中文} & \thead{族} & \thead{она́／оно́／они́} \\
\multicolumn{4}{@{}l}{\scriptsize\color{inkgray}（接上页）7.1　课件 p42 的六组练习} \\
\hline
\endhead
большо́й & 大的 & -ой & {\small больша́я · большо́е · больши́е} \\
хоро́ший & 好的 & жшчщ & {\small хоро́шая · хоро́шее · хоро́шие} \\
све́жий & 新鲜的 & жшчщ & {\small све́жая · све́жее · све́жие} \\
чужо́й & 别人的 & -ой & {\small чужа́я · чужо́е · чужи́е} \\
горя́чий & 热的（物体） & жшчщ & {\small горя́чая · горя́чее · горя́чие} \\
мла́дший & 年小的 & жшчщ & {\small мла́дшая · мла́дшее · мла́дшие} \\
сле́дующий & 下一个的 & жшчщ & {\small сле́дующая · сле́дующее · сле́дующие} \\
\hline
\end{longtable}
\hdr{§8}{УРОК 7（七）：味道与温度}\label{sec:8}
p46 用「甜／酸／苦／咸／辣」五味加冷热去问 \textbf{лимо́н}、\textbf{торт}、\textbf{гру́ша}，答案全是短句。这一组只有 сла́дкий、го́рький 是 г／к／х 族，其余都是硬词干；горя́чий 已在 §7，这里不重复。
\begin{longtable}{@{}p{0.17\linewidth}p{0.19\linewidth}p{0.06\linewidth}p{0.52\linewidth}@{}}
\multicolumn{4}{@{}p{\linewidth}@{}}{\subhdr{8.1　课件 p46 的九个味道词}\par{\small 课件 p46 的模型是：— \textbf{Како́й лимо́н?} — \textbf{Лимо́н жёлтый, ки́слый и невку́сный.} 答案里的形容词按被修饰名词的性数变，所以这一组只要记原形和族。}} \\
\hline
\thead{原形（он）} & \thead{中文} & \thead{族} & \thead{она́／оно́／они́} \\
\hline
\endfirsthead
\hline
\thead{原形（он）} & \thead{中文} & \thead{族} & \thead{она́／оно́／они́} \\
\multicolumn{4}{@{}l}{\scriptsize\color{inkgray}（接上页）8.1　课件 p46 的九个味道词} \\
\hline
\endhead
сла́дкий & 甜的 & гкх & {\small сла́дкая · сла́дкое · сла́дкие} \\
ки́слый & 酸的 & 硬 & {\small ки́слая · ки́слое · ки́слые} \\
го́рький & 苦的 & гкх & {\small го́рькая · го́рькое · го́рькие} \\
солёный & 咸的 & 硬 & {\small солёная · солёное · солёные} \\
о́стрый & 辣的 & 硬 & {\small о́страя · о́строе · о́стрые} \\
холо́дный & 冷的 & 硬 & {\small холо́дная · холо́дное · холо́дные} \\
вку́сный & 好吃的 & 硬 & {\small вку́сная · вку́сное · вку́сные} \\
невку́сный & 不好吃的 & 硬 & {\small невку́сная · невку́сное · невку́сные} \\
\hline
\end{longtable}
\hdr{§9}{УРОК 7（八）：课文与 нра́виться 练习里的形容词}\label{sec:9}
p47 的自我介绍短文与 p48–p51 的 нра́виться 练习把形容词放回句子：\textbf{совреме́нная столи́ца}、\textbf{о́пытный инже́нер}、\textbf{дру́жная гру́ппа}、\textbf{популя́рная му́зыка}、\textbf{пеки́нская о́пера}、\textbf{зи́мняя пого́да}。这一组新词最多，也最接近考试里的组句题。
\begin{longtable}{@{}p{0.17\linewidth}p{0.19\linewidth}p{0.06\linewidth}p{0.52\linewidth}@{}}
\multicolumn{4}{@{}p{\linewidth}@{}}{\subhdr{9.1　课件 p47–p51 的新形容词}\par{\small политехни́ческий 出自 p47 的校名 \textbf{Даля́ньский политехни́ческий университе́т}；词干末尾是 ч，所以走 ж／ш／ч／щ 族。}} \\
\hline
\thead{原形（он）} & \thead{中文} & \thead{族} & \thead{она́／оно́／они́} \\
\hline
\endfirsthead
\hline
\thead{原形（он）} & \thead{中文} & \thead{族} & \thead{она́／оно́／они́} \\
\multicolumn{4}{@{}l}{\scriptsize\color{inkgray}（接上页）9.1　课件 p47–p51 的新形容词} \\
\hline
\endhead
совреме́нный & 现代的 & 硬 & {\small совреме́нная · совреме́нное · совреме́нные} \\
небольшо́й & 不大的 & -ой & {\small небольша́я · небольшо́е · небольши́е} \\
о́пытный & 有经验的 & 硬 & {\small о́пытная · о́пытное · о́пытные} \\
кита́йский & 中国的 & гкх & {\small кита́йская · кита́йское · кита́йские} \\
дру́жный & 友好的、团结的 & жшчщ & {\small дру́жная · дру́жное · дру́жные} \\
популя́рный & 流行的 & 硬 & {\small популя́рная · популя́рное · популя́рные} \\
пеки́нский & 北京的 & гкх & {\small пеки́нская · пеки́нское · пеки́нские} \\
зи́мний & 冬天的 & 软 & {\small зи́мняя · зи́мнее · зи́мние} \\
университе́тский & 大学的 & гкх & {\footnotesize университе́тская · университе́тское · университе́тские} \\
политехни́ческий & 理工的 & жшчщ & {\footnotesize политехни́ческая · политехни́ческое · политехни́ческие} \\
\hline
\end{longtable}
\hdr{§10}{УРОК 8：са́мый 最高级与指示代词里的形容词}\label{sec:10}
У8 把形容词又用了一遍：这一回配 \textbf{са́мый}（са́мый дорого́й／са́мая дешёвая／са́мое тёплое／са́мые хоро́шие）和 \textbf{э́тот／э́та／э́то／э́ти}。新词集中在这一课，共 16 个；颜色、大小、好坏这些旧词不重复列。
\begin{longtable}{@{}p{0.17\linewidth}p{0.19\linewidth}p{0.06\linewidth}p{0.52\linewidth}@{}}
\multicolumn{4}{@{}p{\linewidth}@{}}{\subhdr{10.1　课件 p5–p12、p19–p30 的新形容词}\par{\small са́мый 不是形容词的另一个词尾，而是一个独立的词：它自己按名词的性数变（са́мый／са́мая／са́мое／са́мые），后面的形容词仍写同一个性数——э́тот телефо́н са́мый дорого́й。}} \\
\hline
\thead{原形（он）} & \thead{中文} & \thead{族} & \thead{она́／оно́／они́} \\
\hline
\endfirsthead
\hline
\thead{原形（он）} & \thead{中文} & \thead{族} & \thead{она́／оно́／они́} \\
\multicolumn{4}{@{}l}{\scriptsize\color{inkgray}（接上页）10.1　课件 p5–p12、p19–p30 的新形容词} \\
\hline
\endhead
тёплый & 温暖的 & 硬 & {\small тёплая · тёплое · тёплые} \\
широ́кий & 宽的 & гкх & {\small широ́кая · широ́кое · широ́кие} \\
центра́льный & 中心的 & 硬 & {\small центра́льная · центра́льное · центра́льные} \\
ую́тный & 舒适的、温馨的 & 硬 & {\small ую́тная · ую́тное · ую́тные} \\
дешёвый & 便宜的 & 硬 & {\small дешёвая · дешёвое · дешёвые} \\
шокола́дный & 巧克力的 & 硬 & {\small шокола́дная · шокола́дное · шокола́дные} \\
ко́жаный & 皮革的 & 硬 & {\small ко́жаная · ко́жаное · ко́жаные} \\
ру́сско-кита́йский & 俄汉的 & гкх & {\footnotesize ру́сско-кита́йская · ру́сско-кита́йское · ру́сско-кита́йские} \\
студе́нческий & 学生的 & жшчщ & {\footnotesize студе́нческая · студе́нческое · студе́нческие} \\
люби́мый & 喜爱的 & 硬 & {\small люби́мая · люби́мое · люби́мые} \\
симпати́чный & 好看的、讨喜的 & 硬 & {\small симпати́чная · симпати́чное · симпати́чные} \\
бы́стрый & 快的 & 硬 & {\small бы́страя · бы́строе · бы́стрые} \\
пуши́стый & 毛茸茸的 & 硬 & {\small пуши́стая · пуши́стое · пуши́стые} \\
немо́дный & 不时髦的 & 硬 & {\small немо́дная · немо́дное · немо́дные} \\
удо́бный & 舒适的、方便的 & 硬 & {\small удо́бная · удо́бное · удо́бные} \\
гро́мкий & 响亮的 & гкх & {\small гро́мкая · гро́мкое · гро́мкие} \\
\hline
\end{longtable}
\hdr{§11}{考前速查：五族、拼写、四个坑}\label{sec:11}
这一节不添新词，只把全表的规则压成两页。考前只看这里，加上自己画过道的词，就够了。
\needspace{6\baselineskip}\subhdr{11.1　五族词尾一页表}

\begin{minipage}{\linewidth}\small
\tabtop
\begin{tabular}{@{}p{0.15\linewidth}p{0.18\linewidth}p{0.16\linewidth}p{0.16\linewidth}p{0.16\linewidth}@{}}
\thead{族} & \thead{原形} & \thead{阴性} & \thead{中性} & \thead{复数} \\
\end{tabular}
\tabmid
\begin{tabular}{@{}p{0.15\linewidth}p{0.18\linewidth}p{0.16\linewidth}p{0.16\linewidth}p{0.16\linewidth}@{}}
硬词干 & -ый & -ая & -ое & -ые \\
г／к／х & -ий & -ая & -ое & -ие \\
ж, ш, ч, щ & -ий & -ая & -ее & -ие \\
软词干 & -ий & -яя & -ее & -ие \\
词尾重音 & -о́й & -а́я & -о́е & -ы́е \\
\end{tabular}
\tabbot
\end{minipage}

\textbf{从形式倒推原形}：阴性 -ая／-яя、中性 -ое／-ее、复数 -ые／-ие 把词尾换回 \textbf{-ый／-ий}；-а́я／-о́е／-ы́е 这一族换回 \textbf{-о́й}。例：кра́сная → кра́сный；си́нее → си́ний；голуба́я → голубо́й。
\needspace{6\baselineskip}\subhdr{11.2　拼写规则：七个字母后面不写 ы}

\begin{minipage}{\linewidth}
\tabtop
\begin{tabular}{@{}p{0.34\linewidth}p{0.38\linewidth}p{0.22\linewidth}@{}}
\thead{规则} & \thead{课件例} & \thead{别写成} \\
\end{tabular}
\tabmid
\begin{tabular}{@{}p{0.34\linewidth}p{0.38\linewidth}p{0.22\linewidth}@{}}
词干末尾 г／к／х → -ий & ру́сский, ма́ленький, высо́кий, ти́хий & *ру́скый \\
词干末尾 ж／ш／ч／щ → -ий & хоро́ший, све́жий, горя́чий, мла́дший & *хоро́шый \\
复数也受同一条约束 & больши́е, плохи́е, ти́хие, чужи́е & *большы́е \\
中性写 -ее，不写 -ое & хоро́шее, све́жее, горя́чее & *хоро́шое \\
\end{tabular}
\tabbot
\end{minipage}

\textbf{两个反例}：\textbf{голубы́е} 写 ы——词干末尾是 б，不在七个字母里；\textbf{больши́е} 写 и——词干末尾是 ш。判断只看词干最后一个字母。
\needspace{6\baselineskip}\subhdr{11.3　四个坑（У7 与 У8 都踩到过）}

\begin{itemize}
\item \textbf{坑一：不变格外来词}。\textbf{метро́、пальто́、фо́то、кафе́} 不变格，性按名词规则定：\textbf{Како́е э́то метро́?}、\textbf{Э́то моё пальто́}。
\item \textbf{坑二：ко́фе 是阳性}。虽然以 -е 结尾，它按阳性走：\textbf{Како́й э́то ко́фе?}、\textbf{чёрный ко́фе}。
\item \textbf{坑三：лю́ди 是 челове́к 的复数}。问形容词用复数：\textbf{Каки́е лю́ди?}、\textbf{молоды́е лю́ди}。
\item \textbf{坑四：си́ний 是软词干}。阴性写 \textbf{си́няя}（不是 *си́ная），中性 \textbf{си́нее}，复数 \textbf{си́ние}。
\end{itemize}

\warnbox{\textbf{答题三步}：（1）看问句的疑问词（како́й／кака́я／како́е／каки́е）已经给了性数；（2）选同族词尾；（3）形容词放名词前，不能倒装。构形题（Это … (интересный) книга）先找被修饰的名词，再定族。}
\needspace{6\baselineskip}\subhdr{11.4　са́мый 最高级：它自己也要变性数}

\begin{minipage}{\linewidth}
\tabtop
\begin{tabular}{@{}p{0.12\linewidth}p{0.16\linewidth}p{0.40\linewidth}p{0.24\linewidth}@{}}
\thead{性} & \thead{са́мый} & \thead{课件例句（У8 p2、p5–p8）} & \thead{形容词} \\
\end{tabular}
\tabmid
\begin{tabular}{@{}p{0.12\linewidth}p{0.16\linewidth}p{0.40\linewidth}p{0.24\linewidth}@{}}
阳 & са́мый & Э́тот телефо́н са́мый дорого́й. & дорого́й \\
阴 & са́мая & Э́та арти́стка са́мая краси́вая. & краси́вая \\
中 & са́мое & Э́то пальто́ са́мое тёплое. & тёплое \\
复 & са́мые & Э́ти часы́ са́мые хоро́шие. & хоро́шие \\
\end{tabular}
\tabbot
\end{minipage}

\textbf{一句话}：са́мый 只是「最」这个标记，\textbf{它和后面的形容词必须同性别、同数}——са́мая 配 -ая，са́мое 配 -ое，са́мые 配 -ие。
\hdr{§12}{这份表的来源与边界}\label{sec:12}
形容词的课件分布很集中：У7 p26–p51 是正课，У8 p2–p12、p19–p30 是复现与扩展，前六课只有 добрый、домашний 两个。所以这份表的 §1–§10 就是按这个顺序走的。
\needspace{6\baselineskip}\subhdr{12.1　来源与三条边界}

\begin{itemize}
\item \textbf{来源}：У7（52 页）的 p26–p51 与 У8（30 页）的 p2–p12、p19–p30；每条的原形、族、三个形式都按课件词表与练习抄录，课件没标重音的常用词按标准重音补。
\item \textbf{不补课件外的词}：короткий（短的）、узкий（窄的）、низкий（低的）、скучный（无聊的）这类常用词课件没出现，本表不收，以免把复习范围撑大。
\item \textbf{族的判定}：看\textbf{词干最后一个字母}（去掉词尾 -ый／-ий／-ой 后剩下的部分）与重音位置；重音在词尾的整族换 -о́й，不管词干末尾是什么。
\item \textbf{与名词表配合}：形容词的性数由被修饰的名词决定，名词的性见《俄语名词总表》§2；两表对照看，构形题就不会写错。
\end{itemize}

\aside{这份表按「认形 → 认族 → 认三个形式」设计；三遍读法的具体安排见 §0.3。}
\needspace{6\baselineskip}\subhdr{12.2　十道自测：把括号里的形容词改成需要的形式}

\begin{enumerate}
\item Э́то (но́вый) \blank\ дом.
\item Э́то (но́вый) \blank\ у́лица.
\item Э́то (но́вый) \blank\ окно́.
\item Э́то (но́вый) \blank\ кни́ги.
\item (си́ний) \blank\ не́бо.
\item (большо́й) \blank\ страна́.
\item (хоро́ший) \blank\ друзья́.
\item (ру́сский) \blank\ му́зыка.
\item (голубо́й) \blank\ мо́ре.
\item (ма́ленький) \blank\ го́род.
\end{enumerate}

\warnbox{\textbf{答案}（先自己写，再对）：1 но́вый｜2 но́вая｜3 но́вое｜4 но́вые｜5 си́нее｜6 больша́я｜7 хоро́шие｜8 ру́сская｜9 голубо́е｜10 ма́ленький。\\
\textbf{每题的判断路径}：先看被修饰的名词是什么性、什么数（1–4 分别是阳、阴、中、复），再回到 §0.2 找这个形容词属于哪一族，最后把同族的词尾贴上去。第 10 题是阳性单数，形容词就用原形。}
\end{document}
